\documentclass[12pt, a4paper]{ctexart}
\usepackage{geometry}
\geometry{left=2.5cm, right=2.5cm, top=3cm, bottom=3cm}
\usepackage{listings}
\usepackage{graphicx}
\usepackage{xcolor}

% Code snippet settings
\lstset{
    language=Verilog,
    basicstyle=\ttfamily\small,
    frame=single,
    keywordstyle=\color{blue},
    commentstyle=\color{gray},
    breaklines=true
}

\begin{document}

% 封面
\begin{titlepage}
    \centering
    \vspace*{5cm}
    {\Huge\bfseries 实验报告\par}
    \vspace{2cm}
    {\Large 课程：数字逻辑与计算机组成\par}
    \vspace{1cm}
    {\Large 实验一：多路选择器与优先编码器\par}
    \vspace{4cm}
    {\large 姓名：\underline{\makebox[4cm]{郑袭明}}\par}
    \vspace{1cm}
    {\large 学号：\underline{\makebox[4cm]{251502027}}\par}
    \vspace{1cm}
    {\large 日期：\today\par}
\end{titlepage}

\tableofcontents
\newpage

\section{实验目的}
\begin{enumerate}
    \item 掌握二位四选一选择器的设计与实现。
    \item 掌握 8-3 优先编码器的设计与实现，并结合七段数码管进行结果译码与显示。
    \item 熟悉 FPGA 开发板的拨码开关、LED 和七段数码管的硬件映射与控制。
\end{enumerate}

\section{实验一：二位四选一选择器}
\subsection{设计思路}
利用 \texttt{case} 语句实现 4 选 1 功能。模块包含四个 2 位数据输入 $X0, X1, X2, X3$ 和一个 2 位控制输入 $Y$，根据 $Y$ 的状态选择对应的数据变量输出至 $F$。

\subsection{核心代码}
\begin{lstlisting}
module exp1 (
    input [1:0] X0, X1, X2, X3, Y,
    output reg [1:0] F
);
  always @(X0 or X1 or X2 or X3 or Y)
    case (Y)
      0: F = X0;
      1: F = X1;
      2: F = X2;
      3: F = X3;
      default: F = 2'b0;
    endcase
endmodule
\end{lstlisting}

\subsection{硬件映射}
\begin{itemize}
    \item \textbf{数据输入 ($X0 \sim X3$)}：拨动开关 \texttt{SW[9:2]}
    \item \textbf{控制输入 ($Y$)}：拨动开关 \texttt{SW[1:0]}
    \item \textbf{输出 ($F$)}：\texttt{LED[1:0]}
\end{itemize}

\section{实验二：8-3 优先编码器及数码管显示}
\subsection{设计思路}
实现带使能端（\texttt{en}）的 8-3 高位优先编码器。使用 \texttt{casez} 语句进行含有通配符的模式匹配，当使能有效时输出对应的高位优先编码值 \texttt{value} 及有效指示位 \texttt{valid}。随后利用自定义例化的 \texttt{bcd7seg} 模块将编码值翻译为七段数码管显示控制信号 \texttt{F}。如果输入全0或未使能，则控制显示全灭。

\subsection{核心代码}
\begin{lstlisting}
module enhanced (
    input [7:0] X, input en,
    output reg valid,
    output reg [3:0] value,
    output [6:0] F
);
  always @(X or en) begin
    if (en) begin
      casez (X)
        8'b1???????: begin value = 4'd7; valid = 1'b1; end
        8'b01??????: begin value = 4'd6; valid = 1'b1; end
        8'b001?????: begin value = 4'd5; valid = 1'b1; end
        8'b0001????: begin value = 4'd4; valid = 1'b1; end
        8'b00001???: begin value = 4'd3; valid = 1'b1; end
        8'b000001??: begin value = 4'd2; valid = 1'b1; end
        8'b0000001?: begin value = 4'd1; valid = 1'b1; end
        8'b00000001: begin value = 4'd0; valid = 1'b1; end
        default:     begin value = 4'd0; valid = 1'b0; end
      endcase
    end else begin
      value = 4'd0; valid = 1'b0;
    end
  end

  wire [6:0] F_temp;
  bcd7seg display (.b(value), .h(F_temp));
  assign F = (valid) ? F_temp : 7'b1111111;
endmodule
\end{lstlisting}

\subsection{硬件映射}
\begin{itemize}
    \item \textbf{数据输入 ($X$)}：拨动开关 \texttt{SW[7:0]}
    \item \textbf{使能输入 (\texttt{en})}：拨动开关 \texttt{SW[8]}
    \item \textbf{有效指示 (\texttt{valid})}：\texttt{LED[4]}
    \item \textbf{编码结果输出 (\texttt{value})}：\texttt{LED[3:0]}
    \item \textbf{数码管显示控制}：段选采用译码输出 \texttt{SEG[6:0]}，位选赋固定值 \texttt{AN=8'b11111110} 仅启用最右侧的数码管显示。
\end{itemize}

\end{document}
